% Generated draft: requires mathematical and source review.
\documentclass[11pt]{article}
\usepackage[margin=25mm]{geometry}
\usepackage{fontspec}
\setmainfont{Latin Modern Roman}
\usepackage{amsmath,amssymb,mathtools,unicode-math}
\setmathfont{Latin Modern Math}
\usepackage{enumitem,needspace,xurl,hyperref}
\setlength{\parindent}{0pt}
\setlength{\parskip}{5pt}
\setlength{\emergencystretch}{3em}
\newcommand{\N}{\mathbb{N}}
\newcommand{\Z}{\mathbb{Z}}
\newcommand{\Q}{\mathbb{Q}}
\newcommand{\R}{\mathbb{R}}
\newcommand{\C}{\mathbb{C}}
\newcommand{\F}{\mathbb{F}}
\newcommand{\A}{\mathbb{A}}
\newcommand{\PP}{\mathbb{P}}
\newcommand{\E}{\mathbb{E}}
\newcommand{\eps}{\varepsilon}
\newcommand{\dd}{\,\mathrm{d}}
\begin{document}
\section*{UnsolvedMath ID: 30003315}
\textbf{Source problem code:} catalogue entry
\subsection{Definitions and mathematical statement}
Under the usual convention, a germ at \(+\infty\) of a real-valued function is an equivalence class of functions defined on real intervals \((a,+\infty)\), where two functions are equivalent if they agree for all sufficiently large real arguments. Addition and multiplication of germs are induced by pointwise operations. A Hardy field is a subfield of this ring of germs whose elements have eventually differentiable representatives and which is closed under differentiation.

Let \(F=\bigcup_k F_k\) denote the particular quasianalytic asymptotic field constructed from Ilyashenko algebras in the cited report. The supplied record describes \(F\) as a Hardy field closed under logarithmic composition. Consequently, differentiation is already among its stated closure properties. The record does not supply the definitions of the stages \(F_k\), their indexing range, the construction maps, or the precise meaning of logarithmic composition. Thus this description identifies the intended field by reference to its construction; it does not independently define that field. The germ convention above is standard terminology, but the record does not explicitly specify the germ domain used in the construction.

\textbf{Research question.} Determine further closure properties of this particular field \(F\), beyond the field operations, differentiation, and the logarithmic composition property established in its construction. More explicitly, identify additional operations \(T\) on germs for which, whenever \(f_1,\ldots,f_n\in F\) and \(T(f_1,\ldots,f_n)\) is defined as a germ in the relevant domain, one has
\[
T(f_1,\ldots,f_n)\in F.
\]
The source specifies no particular additional operation or proposed list of operations. This is therefore an open-ended research question, rather than a specified closure conjecture. A fully explicit formulation requires the construction of \(F\) and its precise logarithmic composition convention.
\subsection{Short English statement}
For the specific field \(F\) constructed from Ilyashenko algebras in the cited report, determine which further operations preserve membership in \(F\). Field operations, differentiation, and logarithmic composition are already supplied properties. The construction and the intended additional operations are not specified in the record, so the question remains dependent on the cited source.
\subsection{Sources}
\begin{itemize}
\item Supplied local UnsolvedMath registry record, permanent ID \(30003315\), title \emph{Closure Properties of Logarithmic Hardy Fields}, source problem code \texttt{catalogue entry}, in the collection \emph{Oberwolfach Reports Open Problems}. No individual proposer is supplied.
\item Supplied source citation: \emph{Mini-Workshop: Surreal Numbers, Surreal Analysis, Hahn Fields and Derivations} (2017), Oberwolfach Reports, pp.\ 3313--3372, DOI \texttt{10.4171/owr/2016/60}. Original OWR report: \url{https://doi.org/10.4171/owr/2016/60}. The record gives the proposed year as 2016; the relationship between that year and the citation's 2017 date is not explained.
\item Supplied literature reference, labeled ``Speissegger'': \url{https://doi.org/10.4153/CJM-2016-048-x}. No title or more detailed attribution is supplied for this reference.
\item Literature assessment dated exactly \texttt{2026-08-21}: the supplied curation labels the problem open and states that the construction proves closure under differentiation and logarithmic composition, while the broader request for further closure properties remains open. This is supplied dated curation, not an independently verified literature assessment.
\item The registry labels the statement \texttt{corrected\_verified} and reports that the official passage was checked and rewritten while preserving its content. These are supplied verification claims; no independent source inspection or mathematical verification is asserted here.
\end{itemize}
\end{document}
